\documentclass[12pt,a4paper]{article}
\usepackage[utf8]{inputenc}
\usepackage[T1]{fontenc}
\usepackage[french]{babel}
\usepackage{fontspec}
\usepackage{geometry}
\usepackage{hyperref}
\usepackage{listings}
\usepackage{xcolor}
\usepackage{fontawesome}
\usepackage{setspace}

\geometry{margin=2.5cm,top=3cm,bottom=3cm}
\onehalfspacing

\lstset{
    language=HTML,
    basicstyle=\ttfamily\small,
    keywordstyle=\color{blue},
    commentstyle=\color{green},
    stringstyle=\color{red},
    showstringspaces=false,
    frame=single,
    breaklines=true
}

\title{Activité 1 : Langage HTML - Création d'un site web marchand}
\author{Documentation pour les élèves}
\date{\today}

\begin{document}

\maketitle

\tableofcontents

\newpage

\section{Introduction}

Cette activité vous permet de découvrir les bases du langage HTML en créant un site web marchand simple. Vous partirez d'une structure de pages HTML pré-établies où certains éléments sont à compléter. Votre objectif est de remplacer les textes entre crochets par du contenu approprié et d'ajouter le menu de navigation.

Le site final sera une boutique en ligne avec plusieurs pages interconnectées, une barre de recherche, et un système de suivi des produits consultés via les cookies.

\section{Les balises HTML de base}

Avant de modifier vos pages, il est important de comprendre la structure de base d'un document HTML. Voici les balises essentielles :

\subsection{La balise \texttt{<!DOCTYPE html>}}
Cette déclaration définit le type de document et doit toujours être la première ligne de votre fichier HTML. Elle indique au navigateur qu'il s'agit d'un document HTML5.

\subsection{La balise \texttt{<html>}}
C'est la balise racine qui contient tout le contenu de votre page web. Elle englobe toutes les autres balises.

\subsection{La balise \texttt{<head>}}
Cette section contient les informations "méta" sur votre page :
\begin{itemize}
\item Le titre de la page (\texttt{<title>})
\item Les liens vers les feuilles de style CSS (\texttt{<link>})
\item Les scripts JavaScript (\texttt{<script>})
\item Les informations de codage (\texttt{<meta charset>})
\item Les informations pour les moteurs de recherche (\texttt{<meta>})
\end{itemize}

\textbf{Exemple :}
\begin{lstlisting}
<head>
    <meta charset="UTF-8">
    <meta name="viewport" content="width=device-width, initial-scale=1.0">
    <title>Titre de ma page</title>
    <link rel="stylesheet" href="styles.css">
</head>
\end{lstlisting}

\subsection{La balise \texttt{<body>}}
Cette section contient tout le contenu visible de votre page web :
\begin{itemize}
\item Les titres (\texttt{<h1>}, \texttt{<h2>}, etc.)
\item Les paragraphes (\texttt{<p>})
\item Les images (\texttt{<img>})
\item Les liens (\texttt{<a>})
\item Les formulaires (\texttt{<form>})
\item Les listes (\texttt{<ul>}, \texttt{<ol>})
\item Les divisions (\texttt{<div>})
\end{itemize}

\textbf{Exemple de structure complète :}
\begin{lstlisting}
<!DOCTYPE html>
<html lang="fr">
<head>
    <meta charset="UTF-8">
    <title>Ma première page</title>
</head>
<body>
    <h1>Mon titre principal</h1>
    <p>Mon premier paragraphe.</p>
</body>
</html>
\end{lstlisting}

\subsection{Conseils pour bien débuter}
\begin{itemize}
\item Respectez toujours l'ordre des balises : \texttt{<!DOCTYPE html>} → \texttt{<html>} → \texttt{<head>} → \texttt{<body>}
\item Fermez toujours vos balises avec la barre oblique : \texttt{</balise>}
\item Utilisez l'indentation pour rendre votre code plus lisible
\item Testez régulièrement vos modifications dans un navigateur
\end{itemize}

\section{Structure générale du site}

\subsection{Arborescence du site}

Le site est organisé selon la structure suivante :

\begin{itemize}
\item \textbf{index.html} : Page d'accueil
\item \textbf{pages/}
  \begin{itemize}
  \item inscription.html : Formulaire d'inscription
  \item panier.html : Gestion du panier
  \item plan.html : Plan du site
  \item produit\_detail.html : Détail d'un produit
  \item produits.html : Catalogue des produits
  \item recherche.html : Résultats de recherche
  \end{itemize}
\end{itemize}

\subsection{Éléments communs à toutes les pages}

Chaque page contient plusieurs sections communes :

\subsubsection{Barre de recherche}
Située dans l'en-tête, elle permet aux utilisateurs de rechercher des produits. Elle contient :
\begin{itemize}
\item Un champ de saisie (\texttt{<input type="text">})
\item Un bouton "Rechercher"
\item Un bouton "Voir l'index"
\end{itemize}

\subsubsection{Section des produits récemment consultés}
Cette section utilise les cookies pour mémoriser les produits visités par l'utilisateur. Elle comprend :
\begin{itemize}
\item Un titre explicatif
\item Des boutons pour afficher ou effacer les cookies
\item Une grille où s'affichent les produits récents
\end{itemize}

\subsubsection{Menu de navigation}

Le menu de navigation est un élément essentiel de votre site web. Il permet aux utilisateurs de naviguer facilement entre les différentes pages. Le menu doit être ajouté à l'emplacement indiqué par le commentaire \texttt{<!-- AJOUTER LE MENU ICI -->} dans chaque page.

\textbf{Structure du menu :}
\begin{itemize}
\item Utilisez la balise \texttt{<nav>} avec l'id \texttt{main-nav}
\item Créez des liens \texttt{<a>} vers chaque page principale
\item Respectez la hiérarchie des chemins relatifs selon la page où vous vous trouvez
\end{itemize}

\textbf{Exemple de menu pour la page d'accueil (index.html) :}

\begin{lstlisting}
<nav id="main-nav">
    <a href="index.html">Accueil</a>
    <a href="pages/produits.html">Produits</a>
    <a href="pages/panier.html">Panier \faShoppingCart{}</a>
    <a href="pages/inscription.html">Inscription</a>
    <a href="pages/plan.html">Plan du site</a>
</nav>
\end{lstlisting}

\textbf{Exemple de menu pour les pages dans le dossier pages/ :}

\begin{lstlisting}
<nav id="main-nav">
    <a href="../index.html">Accueil</a>
    <a href="produits.html">Produits</a>
    <a href="panier.html">Panier \faShoppingCart{}</a>
    <a href="inscription.html">Inscription</a>
    <a href="plan.html">Plan du site</a>
</nav>
\end{lstlisting}

\textbf{À propos des chemins relatifs :}
\begin{itemize}
\item \texttt{../} permet de remonter d'un niveau dans l'arborescence
\item Les chemins doivent être adaptés selon l'emplacement de chaque page
\item Testez toujours vos liens en ouvrant les pages dans un navigateur
\end{itemize}

\section{Rôle des différentes pages}

\subsection{Page d'accueil (index.html)}
C'est la page principale du site. Elle présente :
\begin{itemize}
\item Le titre de la boutique
\item Une section "héros" avec un message de bienvenue
\item Une section de présentation
\item Des boutons d'action (vers les produits et l'inscription)
\item La section des produits récemment consultés
\end{itemize}

\subsection{Page d'inscription (pages/inscription.html)}
Permet aux visiteurs de s'inscrire en tant que clients. Contient :
\begin{itemize}
\item Un formulaire avec champs pour prénom, nom, email, ville et âge
\item Un bouton de soumission
\item Une section affichant la liste des inscrits
\item La section des produits récemment consultés
\end{itemize}

\subsection{Page panier (pages/panier.html)}
Affiche le contenu du panier de l'utilisateur :
\begin{itemize}
\item La liste des articles sélectionnés
\item Le total du panier
\item Des boutons pour vider ou valider la commande
\item La section des produits récemment consultés
\end{itemize}

\subsection{Page plan du site (pages/plan.html)}
Présente la structure du site :
\begin{itemize}
\item Liste des pages principales avec descriptions
\item Informations sur le projet pédagogique
\item Liste des fonctionnalités
\item La section des produits récemment consultés
\end{itemize}


\section{Ajouter des images à vos pages}

Pour enrichir vos pages web avec des images, vous pouvez utiliser la balise \texttt{<img>}. Voici comment procéder :

\textbf{Syntaxe de base :}
\begin{lstlisting}
<img src="chemin/vers/image.jpg" alt="Description de l'image">
\end{lstlisting}

\textbf{Attributs importants :}
\begin{itemize}
\item \texttt{src} : chemin vers le fichier image (relatif ou absolu)
\item \texttt{alt} : texte alternatif décrivant l'image (obligatoire pour l'accessibilité)
\item \texttt{width} et \texttt{height} : dimensions de l'image (optionnel)
\end{itemize}

\textbf{Exemples d'utilisation :}
\begin{itemize}
\item Image dans le même dossier : \texttt{<img src="logo.png" alt="Logo de la boutique">}
\item Image dans un sous-dossier : \texttt{<img src="images/photo.jpg" alt="Photo du produit">}
\item Image dans le dossier parent : \texttt{<img src="../images/logo.png" alt="Logo principal">}
\end{itemize}

\textbf{Formats d'images supportés :}
\begin{itemize}
\item JPEG (.jpg, .jpeg) : pour les photos
\item PNG (.png) : pour les images avec transparence
\item GIF (.gif) : pour les animations simples
\item WebP (.webp) : format moderne optimisé
\end{itemize}

\textbf{Conseils pour les images :}
\begin{itemize}
\item Utilisez des images de taille appropriée pour éviter de ralentir le chargement
\item Toujours spécifier l'attribut \texttt{alt} pour l'accessibilité
\item Organisez vos images dans un dossier dédié (ex: \texttt{images/})
\end{itemize}

\section{Instructions pour modifier les pages}

Pour chaque page, vous devez remplacer les textes entre crochets \texttt{[...]} par du contenu approprié en français. Voici le détail pour chaque fichier :

\subsubsection{index.html}

\textbf{Objectif :} Cette page constitue la vitrine de votre boutique en ligne. Elle doit accueillir les visiteurs et les inciter à explorer le catalogue de produits.

\begin{itemize}
\item \texttt{[TITRE DE LA PAGE]} : "Boutique SNT - Accueil" \newline \textit{Objectif : Définir le titre qui apparaîtra dans l'onglet du navigateur}
\item \texttt{[TITRE PRINCIPAL]} : "\faShoppingCart{} Boutique SNT" \newline \textit{Objectif : Donner un nom accrocheur à votre boutique avec un symbole représentatif}
\item \texttt{[TITRE DE LA SECTION HERO]} : "Bienvenue sur notre boutique en ligne !" \newline \textit{Objectif : Accueillir chaleureusement les visiteurs dès leur arrivée}
\item \texttt{[TEXTE DE PRESENTATION]} : "Découvrez nos produits et faites vos achats en toute simplicité." \newline \textit{Objectif : Présenter brièvement les avantages de votre boutique}
\item \texttt{[TITRE DE LA SECTION PRESENTATION]} : "À propos de notre boutique" \newline \textit{Objectif : Introduire la section descriptive de votre commerce}
\item \texttt{[PARAGRAPHE 1]} : "Notre boutique propose une sélection de produits de qualité." \newline \textit{Objectif : Mettre en avant la qualité de vos produits}
\item \texttt{[PARAGRAPHE 2]} : "Naviguez dans notre catalogue et ajoutez vos articles préférés au panier !" \newline \textit{Objectif : Guider l'utilisateur vers l'action d'achat}
\item \texttt{[TEXTE BOUTON 1]} : "Voir les produits" \newline \textit{Objectif : Inciter à découvrir le catalogue}
\item \texttt{[TEXTE BOUTON 2]} : "S'inscrire" \newline \textit{Objectif : Encourager l'inscription des nouveaux clients}
\item \texttt{[TITRE SECTION HISTORIQUE]} : "\faClockO{} Vous avez récemment consulté" \newline \textit{Objectif : Rappeler aux utilisateurs leurs consultations récentes}
\item \texttt{[TEXTE INFO COOKIES]} : "Ces produits sont mémorisés dans un cookie sur votre navigateur" \newline \textit{Objectif : Expliquer le fonctionnement technique des cookies}
\item \texttt{[TEXTE BOUTON COOKIES]} : "Voir mes cookies 🍪" \newline \textit{Objectif : Permettre la visualisation des données stockées}
\item \texttt{[TEXTE BOUTON EFFACER]} : "Effacer l'historique" \newline \textit{Objectif : Offrir la possibilité de supprimer les données personnelles}
\item \texttt{[TEXTE DU FOOTER]} : "© 2025 Boutique SNT - Projet pédagogique" \newline \textit{Objectif : Afficher les droits d'auteur et la nature éducative du projet}
\end{itemize}

Ajoutez le menu de navigation après le commentaire \texttt{<!-- AJOUTER LE MENU ICI -->} :

\begin{lstlisting}
<nav id="main-nav">
    <a href="index.html">Accueil</a>
    <a href="pages/produits.html">Produits</a>
    <a href="pages/panier.html">Panier \faShoppingCart{}</a>
    <a href="pages/inscription.html">Inscription</a>
    <a href="pages/plan.html">Plan du site</a>
</nav>
\end{lstlisting}

\subsubsection{pages/inscription.html}

\textbf{Objectif :} Cette page permet aux visiteurs de créer un compte client en fournissant leurs informations personnelles.

\begin{itemize}
\item \texttt{[TITRE DE LA PAGE]} : "Boutique SNT - Inscription" \newline \textit{Objectif : Indiquer clairement la fonction de la page dans l'onglet}
\item \texttt{[TITRE PRINCIPAL]} : "🛒 Boutique SNT" \newline \textit{Objectif : Maintenir la cohérence de la marque sur toutes les pages}
\item \texttt{[TITRE PRINCIPAL DE LA PAGE]} : "Inscription" \newline \textit{Objectif : Annoncer clairement l'action disponible sur cette page}
\item \texttt{[LABEL PRENOM]} : "Prénom :" \newline \textit{Objectif : Guider l'utilisateur pour saisir son prénom}
\item \texttt{[LABEL NOM]} : "Nom :" \newline \textit{Objectif : Demander le nom de famille de l'utilisateur}
\item \texttt{[LABEL EMAIL]} : "Email :" \newline \textit{Objectif : Collecter l'adresse email pour les communications}
\item \texttt{[LABEL VILLE]} : "Ville :" \newline \textit{Objectif : Connaître la localisation géographique du client}
\item \texttt{[LABEL AGE]} : "Âge :" \newline \textit{Objectif : Vérifier l'éligibilité selon les critères d'âge}
\item \texttt{[TEXTE BOUTON]} : "S'inscrire" \newline \textit{Objectif : Valider et soumettre le formulaire d'inscription}
\item \texttt{[TITRE LISTE INSCRITS]} : "Liste des inscrits" \newline \textit{Objectif : Présenter les clients déjà enregistrés}
\item \texttt{[TITRE SECTION HISTORIQUE]} : "🕒 Vous avez récemment consulté" \newline \textit{Objectif : Montrer l'historique de navigation personnel}
\item \texttt{[TEXTE INFO COOKIES]} : "Ces produits sont mémorisés dans un cookie sur votre navigateur" \newline \textit{Objectif : Expliquer la persistance des données}
\item \texttt{[TEXTE BOUTON COOKIES]} : "Voir mes cookies 🍪" \newline \textit{Objectif : Permettre l'inspection des données stockées}
\item \texttt{[TEXTE BOUTON EFFACER]} : "Effacer l'historique" \newline \textit{Objectif : Respecter la confidentialité des données}
\item \texttt{[TEXTE DU FOOTER]} : "© 2025 Boutique SNT - Projet pédagogique" \newline \textit{Objectif : Affirmer les droits et la nature éducative}
\end{itemize}

Ajoutez le même menu de navigation qu'à l'accueil, en ajustant les chemins relatifs si nécessaire.

\subsubsection{pages/panier.html}

\textbf{Objectif :} Cette page affiche le contenu actuel du panier d'achat et permet de gérer les articles sélectionnés.

\begin{itemize}
\item \texttt{[TITRE DE LA PAGE]} : "Boutique SNT - Panier" \newline \textit{Objectif : Indiquer la fonction de gestion du panier}
\item \texttt{[TITRE PRINCIPAL]} : "🛒 Boutique SNT" \newline \textit{Objectif : Maintenir l'identité visuelle de la boutique}
\item \texttt{[TITRE PRINCIPAL DE LA PAGE]} : "Mon Panier" \newline \textit{Objectif : Personnaliser l'expérience utilisateur}
\item \texttt{[TEXTE TOTAL]} : "Total :" \newline \textit{Objectif : Présenter le montant total de la commande}
\item \texttt{[TEXTE BOUTON VIDER]} : "Vider le panier" \newline \textit{Objectif : Permettre de supprimer tous les articles}
\item \texttt{[TEXTE BOUTON VALIDER]} : "Valider la commande" \newline \textit{Objectif : Finaliser le processus d'achat}
\item \texttt{[TITRE SECTION HISTORIQUE]} : "🕒 Vous avez récemment consulté" \newline \textit{Objectif : Rappeler les produits vus récemment}
\item \texttt{[TEXTE INFO COOKIES]} : "Ces produits sont mémorisés dans un cookie sur votre navigateur" \newline \textit{Objectif : Expliquer le mécanisme de stockage}
\item \texttt{[TEXTE BOUTON COOKIES]} : "Voir mes cookies 🍪" \newline \textit{Objectif : Offrir la transparence sur les données}
\item \texttt{[TEXTE BOUTON EFFACER]} : "Effacer l'historique" \newline \textit{Objectif : Respecter les droits de suppression des données}
\item \texttt{[TEXTE DU FOOTER]} : "© 2025 Boutique SNT - Projet pédagogique" \newline \textit{Objectif : Conserver la signature du projet}
\end{itemize}

Ajoutez le menu de navigation.

\subsubsection{pages/plan.html}

\textbf{Objectif :} Cette page présente la structure complète du site web et explique ses fonctionnalités pédagogiques.

\begin{itemize}
\item \texttt{[TITRE DE LA PAGE]} : "Boutique SNT - Plan du site" \newline \textit{Objectif : Définir le rôle de navigation de cette page}
\item \texttt{[TITRE PRINCIPAL]} : "🛒 Boutique SNT" \newline \textit{Objectif : Assurer la cohérence de l'en-tête}
\item \texttt{[TITRE PRINCIPAL DE LA PAGE]} : "Plan du Site" \newline \textit{Objectif : Annoncer clairement le contenu informatif}
\item \texttt{[TITRE SECTION 1]} : "\faMobile{} Pages principales" \newline \textit{Objectif : Organiser la présentation des pages du site}
\item \texttt{[LIEN ACCUEIL]} : "\faHome{} Accueil - Page de bienvenue" \newline \textit{Objectif : Décrire la fonction de la page d'accueil}
\item \texttt{[LIEN PRODUITS]} : "\faArchive{} Produits - Catalogue des produits disponibles" \newline \textit{Objectif : Présenter le catalogue comme source d'information}
\item \texttt{[LIEN PANIER]} : "\faShoppingCart{} Panier - Vos articles sélectionnés" \newline \textit{Objectif : Expliquer l'utilité du panier d'achat}
\item \texttt{[LIEN INSCRIPTION]} : "\faPencil{} Inscription - Créer un compte client" \newline \textit{Objectif : Décrire le processus de création de compte}
\item \texttt{[LIEN PLAN]} : "\faMap{} Plan du site - Cette page" \newline \textit{Objectif : S'auto-référencer pour la clarté}
\item \texttt{[TITRE SECTION 2]} : "\faInfo{} Informations" \newline \textit{Objectif : Introduire les informations pédagogiques}
\item \texttt{[PARAGRAPHE 1]} : "Ce site est un projet pédagogique réalisé dans le cadre du cours de Sciences Numériques et Technologiques." \newline \textit{Objectif : Situer le contexte éducatif du projet}
\item \texttt{[PARAGRAPHE 2]} : "Il permet d'apprendre les bases du développement web (HTML, CSS, JavaScript) et de la programmation backend (Python)." \newline \textit{Objectif : Lister les technologies abordées}
\item \texttt{[TITRE SECTION 3]} : "\faBullseye{} Fonctionnalités" \newline \textit{Objectif : Présenter les capacités techniques du site}
\item \texttt{[FONCTIONNALITE 1]} : "Affichage dynamique des produits depuis un fichier YAML" \newline \textit{Objectif : Montrer l'utilisation de données externes}
\item \texttt{[FONCTIONNALITE 2]} : "Ajout d'articles au panier" \newline \textit{Objectif : Décrire l'interaction utilisateur de base}
\item \texttt{[FONCTIONNALITE 3]} : "Gestion du panier (ajout, suppression, calcul du total)" \newline \textit{Objectif : Expliquer la logique métier}
\item \texttt{[FONCTIONNALITE 4]} : "Inscription de nouveaux clients" \newline \textit{Objectif : Présenter la gestion des utilisateurs}
\item \texttt{[FONCTIONNALITE 5]} : "Stockage des données clients dans un fichier YAML" \newline \textit{Objectif : Montrer la persistance des données}
\item \texttt{[TITRE SECTION HISTORIQUE]} : "🕒 Vous avez récemment consulté" \newline \textit{Objectif : Maintenir la cohérence avec les autres pages}
\item \texttt{[TEXTE INFO COOKIES]} : "Ces produits sont mémorisés dans un cookie sur votre navigateur" \newline \textit{Objectif : Expliquer le mécanisme technique}
\item \texttt{[TEXTE BOUTON COOKIES]} : "Voir mes cookies 🍪" \newline \textit{Objectif : Permettre l'inspection des données}
\item \texttt{[TEXTE BOUTON EFFACER]} : "Effacer l'historique" \newline \textit{Objectif : Respecter la protection des données}
\item \texttt{[TEXTE DU FOOTER]} : "© 2025 Boutique SNT - Projet pédagogique" \newline \textit{Objectif : Affirmer l'origine éducative}
\end{itemize}

Ajoutez le menu de navigation.

\subsection{pages/produit\_detail.html}

\begin{itemize}
\item \texttt{[TITRE DE LA PAGE]} : "Boutique SNT - Détail du produit"
\item \texttt{[TITRE PRINCIPAL]} : "🛒 Boutique SNT"
\item \texttt{[TEXTE CHARGEMENT]} : "Chargement..."
\item \texttt{[TITRE SECTION HISTORIQUE]} : "🕒 Vous avez récemment consulté"
\item \texttt{[TEXTE INFO COOKIES]} : "Ces produits sont mémorisés dans un cookie sur votre navigateur"
\item \texttt{[TEXTE BOUTON COOKIES]} : "Voir mes cookies 🍪"
\item \texttt{[TEXTE BOUTON EFFACER]} : "Effacer l'historique"
\item \texttt{[TEXTE DU FOOTER]} : "© 2025 Boutique SNT - Projet pédagogique"
\end{itemize}

Ajoutez le menu de navigation.

\subsection{pages/produits.html}

\begin{itemize}
\item \texttt{[TITRE DE LA PAGE]} : "Boutique SNT - Produits"
\item \texttt{[TITRE PRINCIPAL]} : "🛒 Boutique SNT"
\item \texttt{[TITRE PRINCIPAL DE LA PAGE]} : "Nos Produits"
\item \texttt{[TITRE SECTION HISTORIQUE]} : "🕒 Vous avez récemment consulté"
\item \texttt{[TEXTE INFO COOKIES]} : "Ces produits sont mémorisés dans un cookie sur votre navigateur"
\item \texttt{[TEXTE BOUTON COOKIES]} : "Voir mes cookies 🍪"
\item \texttt{[TEXTE BOUTON EFFACER]} : "Effacer l'historique"
\item \texttt{[TEXTE DU FOOTER]} : "© 2025 Boutique SNT - Projet pédagogique"
\end{itemize}

Ajoutez le menu de navigation.

\subsection{pages/recherche.html}

\begin{itemize}
\item \texttt{[TITRE DE LA PAGE]} : "Boutique SNT - Résultats de recherche"
\item \texttt{[TITRE PRINCIPAL]} : "🛒 Boutique SNT"
\item \texttt{[TITRE PRINCIPAL DE LA PAGE]} : "Résultats de recherche"
\item \texttt{[TITRE SECTION HISTORIQUE]} : "🕒 Vous avez récemment consulté"
\item \texttt{[TEXTE INFO COOKIES]} : "Ces produits sont mémorisés dans un cookie sur votre navigateur"
\item \texttt{[TEXTE BOUTON COOKIES]} : "Voir mes cookies 🍪"
\item \texttt{[TEXTE BOUTON EFFACER]} : "Effacer l'historique"
\item \texttt{[TEXTE DU FOOTER]} : "© 2025 Boutique SNT - Projet pédagogique"
\end{itemize}

Ajoutez le menu de navigation.

\end{document}